\documentclass[11pt, a4paper]{article}

% --- Packages ---
\usepackage[margin=2.5cm]{geometry}
\usepackage{amsmath}
\usepackage{amssymb}
\usepackage{booktabs}
\usepackage{array}
\usepackage{parskip}
\usepackage{hyperref}
\usepackage{microtype}

% --- Hyperref setup ---
\hypersetup{
    colorlinks=true,
    linkcolor=black,
    urlcolor=blue,
    citecolor=black
}

% --- Title ---
\title{\textbf{Quantitative Dynamic Macroeconomics}\\[0.4em]
\large Baseline Model Handout}
\author{Konstantin Gantert}
\date{this version: April 2026}

% ---------------------------------------------------------------
\begin{document}
\maketitle

This handout describes the baseline New Keynesian DSGE model that serves as the
starting point for all project topics. It presents the optimisation problems,
constraints, and first-order conditions (FOCs) of each agent, followed by the
general equilibrium conditions and policy rules. Each project topic extends this
baseline by modifying or adding exactly one block. The notation here is the
reference for all five topics.

An important property: \textbf{all topic extensions are designed so that
the steady state of the extended model is identical to that of the baseline}.
Only the dynamics of the model change. Students can use this as a verification
check --- if your extended model produces a different steady state, something
has gone wrong in the implementation.

% ---------------------------------------------------------------
\section{Households}

The representative household maximises lifetime utility over consumption $C_t$
and leisure, subject to a cash-in-advance (CIA) constraint, a budget constraint,
and a capital accumulation equation with quadratic investment adjustment costs:
%
\begin{equation}
    \max_{\{C_t,\,L_t,\,X_t,\,M_{t+1}\}} \;
    E_0 \sum_{t=0}^{\infty} \beta^t
    \left[
        \frac{C_t^{1-\sigma}-1}{1-\sigma}
        -
        \frac{\phi}{1+\gamma} L_t^{1+\gamma}
    \right]
\end{equation}
%
subject to the \textbf{cash-in-advance constraint:}
\begin{equation}
    \frac{M_t}{\nu} \;\geq\; P_t(1+\tau_C)\,C_t
    \label{eq:CIA}
\end{equation}
%
the \textbf{budget constraint:}
\begin{equation}
    (1-\tau_L)W_t L_t
    + (1-\tau_K)P_t r_t K_t
    + M_t + D_t - \tau_t
    \;\geq\;
    M_{t+1}
    + P_t(1+\tau_C)\,C_t
    + P_t X_t
    \label{eq:budget}
\end{equation}
%
and the \textbf{capital law of motion with investment adjustment costs:}
\begin{equation}
    K_{t+1} = (1-\delta)\,K_t
    + \left(1 - \frac{\varphi}{2}
      \left(\frac{X_t}{X_{t-1}}-1\right)^{\!2}\right) X_t
    \label{eq:capital}
\end{equation}
%
where $\varphi \geq 0$ governs the magnitude of investment adjustment costs.
At $\varphi = 0$ the standard capital law of motion without adjustment costs
is recovered. At $\varphi > 0$ it is costly to change investment rapidly,
making the capital stock respond sluggishly to shocks. The lagged investment
level $X_{t-1}$ is a predetermined state variable in the baseline.

\subsection*{Notation}

$C_t$: consumption; $L_t$: hours worked; $M_t$: money holdings at the beginning
of period $t$; $P_t$: price level; $W_t$: nominal wage; $r_t$: real rental rate
of capital services; $K_t$: capital stock at the beginning of period $t$;
$X_t$: capital investment; $D_t$: dividends; $\beta \in (0,1)$: discount
factor; $\sigma > 0$: coefficient of relative risk aversion; $\phi > 0$: labor
disutility weight; $\gamma > 0$: inverse Frisch elasticity of labor supply;
$\nu > 0$: money-to-consumption ratio (inverse money velocity); $\delta \in (0,1)$: capital depreciation rate;
$\tau_C,\,\tau_L,\,\tau_K \geq 0$: fixed tax rates on consumption, labor
income, and capital income respectively.

\subsection*{First-Order Conditions}

Let $\tilde{\mu}_t \geq 0$ denote the Lagrange multiplier on the CIA constraint
and $\tilde{\lambda}_t \geq 0$ the multiplier on the budget constraint, both in
stationary form. Let $\tilde{\Theta}_t$ denote the Lagrange multiplier on the
capital law of motion (\ref{eq:capital}). Let $w_t = W_t/P_t$ be the real wage.

\textbf{CIA and budget multipliers:}
\begin{equation}
    1 = \tilde{\mu}_t + \tilde{\lambda}_t
    \label{eq:mult_split}
\end{equation}

\textbf{Labor supply:}
\begin{equation}
    \tilde{\mu}_t
    = \frac{(1+\tau_C)\,\phi\,L_t^{\gamma}\,C_t^{\sigma}}{(1-\tau_L)\,w_t}
    \label{eq:labor_supply}
\end{equation}

\textbf{Money demand (Euler equation for money):}
\begin{equation}
    \tilde{\mu}_t
    = \beta\,E_t\!\left[
        \left(\tilde{\mu}_{t+1} + \frac{\tilde{\lambda}_{t+1}}{\nu}\right)
        \frac{P_t}{P_{t+1}}
        \left(\frac{C_t}{C_{t+1}}\right)^{\!\sigma}
    \right]
    \label{eq:money_euler}
\end{equation}

\textbf{Bond pricing (Euler equation for bonds):}
\begin{equation}
    \frac{1}{1+r_{B,t}}
    = \beta\,E_t\!\left[
        \frac{\tilde{\mu}_{t+1}}{\tilde{\mu}_t}
        \frac{P_t}{P_{t+1}}
        \left(\frac{C_t}{C_{t+1}}\right)^{\!\sigma}
    \right]
    \label{eq:bond_euler}
\end{equation}

\textbf{Investment (FOC for $X_t$) --- Tobin's $Q$:}
\begin{align}
    \tilde{\mu}_t
    &= \tilde{\Theta}_t
       \left(1 - \frac{\varphi}{2}\left(\frac{X_t}{X_{t-1}}-1\right)^{\!2}
             - \varphi\left(\frac{X_t}{X_{t-1}}-1\right)\frac{X_t}{X_{t-1}}
       \right) \notag \\
    &\quad
    + \beta\,E_t\!\left[
        \left(\frac{C_t}{C_{t+1}}\right)^{\!\sigma}
        \tilde{\Theta}_{t+1}\,\varphi
        \left(\frac{X_{t+1}}{X_t}-1\right)
        \left(\frac{X_{t+1}}{X_t}\right)^{\!2}
    \right]
    \label{eq:investment_foc}
\end{align}

\textbf{Capital (Euler equation for $K_{t+1}$):}
\begin{equation}
    \tilde{\Theta}_t
    = \beta\,E_t\!\left[
        \left(\frac{C_t}{C_{t+1}}\right)^{\!\sigma}
        \left((1-\delta)\tilde{\Theta}_{t+1} + (1-\tau_K)\,r_{t+1}\tilde{\mu}_{t+1}\right)
    \right]
    \label{eq:capital_euler}
\end{equation}

\textbf{Stochastic discount factor:}
\begin{equation}
    \beta_{t,t+1} = \frac{1}{1+r_{B,t}}
    \label{eq:sdf}
\end{equation}

\subsection*{Tobin's $Q$}

The Lagrange multiplier $\tilde{\Theta}_t$ on the capital law of motion
(\ref{eq:capital}) is the shadow price of an additional unit of installed
capital --- this is \textbf{Tobin's $Q$}. It represents the ratio of the market
value of installed capital to the cost of a new investment good. When
$\tilde{\Theta}_t > \tilde{\mu}_t$, installed capital is worth more than new
investment goods, providing an incentive to invest. At steady state
$\tilde{\Theta} = \tilde{\mu}$, so Tobin's $Q$ equals one and there is no
incentive to invest beyond replacement of depreciated capital. Investment
adjustment costs make Tobin's $Q$ a genuine state-like variable --- its response
to shocks is hump-shaped rather than immediate, because rapid investment changes
are costly. This is why IRFs for investment under IAC display sluggish, persistent
dynamics that more closely match empirical business cycle patterns than the
frictionless model.

% ---------------------------------------------------------------
\section{Final Goods Firm}

A representative competitive firm aggregates a continuum of differentiated
intermediate goods into a homogeneous final good via a CES technology:
%
\begin{equation}
    Y_t = \left(\int_0^1 Y_t(i)^{\rho}\,di\right)^{1/\rho}
\end{equation}
%
where $\rho \in (0,1)$ is the inverse of the steady-state gross price markup.
Profit maximisation yields the \textbf{demand schedule} for variety $i$:
%
\begin{equation}
    \frac{P_t(i)}{P_t}
    = \left(\frac{Y_t}{Y_t(i)}\right)^{1-\rho}
    \label{eq:demand}
\end{equation}

% ---------------------------------------------------------------
\section{Intermediate Goods Firms}

A continuum of monopolistically competitive firms $i \in [0,1]$ each produce a
differentiated variety. Firm $i$ maximises the present discounted value of
profits:
%
\begin{equation}
    \max_{\{Y_t(i),\,L_t(i),\,K_t(i),\,P_t(i)\}}
    \;E_0\sum_{t=0}^{\infty} \beta_{0,t}
    \Bigl[P_t(i)\,Y_t(i) - W_t\,L_t(i) - P_t\,r_t\,K_t(i)\Bigr]
\end{equation}
%
subject to the \textbf{production technology with Rotemberg price adjustment
costs:}
\begin{equation}
    Y_t(i) = \bigl(Z_t\,L_t(i)\bigr)^{1-\alpha} K_t(i)^{\alpha}
    \left[1 - \frac{\kappa}{2}
    \left(\frac{P_t(i)}{P_{t-1}(i)} - 1\right)^{\!2}\right]
    \label{eq:production}
\end{equation}
%
and the \textbf{demand constraint:}
\begin{equation}
    Y_t(i) = \left(\frac{P_t(i)}{P_t}\right)^{\!\frac{1}{\rho-1}} Y_t
    \label{eq:demand_constraint}
\end{equation}

\subsection*{Notation}

$Z_t$: total factor productivity; $\alpha \in (0,1)$: capital elasticity of
output; $\kappa \geq 0$: price adjustment cost parameter; $mc_t$: real marginal
cost; $\pi_t$: net inflation rate defined by $1+\pi_t = P_t/P_{t-1}$.

\subsection*{First-Order Conditions}

Imposing symmetry ($P_t(i)=P_t$, $Y_t(i)=Y_t$, $L_t(i)=L_t$, $K_t(i)=K_t$):

\textbf{Labor demand:}
\begin{equation}
    w_t = (1-\alpha)\,\frac{Y_t}{L_t}\,mc_t
    \label{eq:labor_demand}
\end{equation}

\textbf{Capital demand:}
\begin{equation}
    r_t = \alpha\,\frac{Y_t}{K_t}\,mc_t
    \label{eq:capital_demand}
\end{equation}

\textbf{New Keynesian Phillips curve (Rotemberg):}
\begin{equation}
    \frac{\kappa\,\pi_t\,(1+\pi_t)}{1-\frac{\kappa}{2}\pi_t^2}
    =
    \frac{1 - \rho / mc_t}{1-\rho}
    +
    E_t\!\left[
        \beta_{t,t+1}\,(1+\pi_{t+1})\,\frac{Y_{t+1} \, mc_{t+1}}{Y_t \, mc_t}\,
        \frac{\kappa\,\pi_{t+1}\,(1+\pi_{t+1})}{1-\frac{\kappa}{2}\pi_{t+1}^2}
    \right]
    \label{eq:phillips}
\end{equation}

% ---------------------------------------------------------------
\section{General Equilibrium}

\textbf{Resource constraint:}
\begin{equation}
    C_t + X_t = Y_t
    \label{eq:resource}
\end{equation}

\textbf{Fiscal policy} (balanced budget, no government debt):
\begin{equation}
    \tau_t = \tau_C C_t + \tau_L w_t L_t + \tau_K r_t K_t
    \label{eq:fiscal}
\end{equation}

\textbf{Money market} (CIA constraint binds in equilibrium):
\begin{equation}
    \frac{M_t}{\nu} = P_t(1+\tau_C)\,C_t
    \label{eq:money_market}
\end{equation}

\textbf{Inflation identity:}
\begin{equation}
    1 + \pi_t = \frac{P_t}{P_{t-1}}
    \label{eq:inflation}
\end{equation}

% ---------------------------------------------------------------
\section{Policy Rules and Exogenous Processes}

\textbf{Monetary policy} (Taylor-type interest rate rule):
\begin{equation}
    \frac{1+r_{B,t}}{1+\bar{r}_B}
    =
    \left(\frac{1+r_{B,t-1}}{1+\bar{r}_B}\right)^{\!\theta_R}
    \!\left[
        \left(\frac{1+\pi_t}{1+\pi^*}\right)^{\!\theta_\pi}
        \left(\frac{Y_t}{\bar{Y}}\right)^{\!\theta_Y}
    \right]^{1-\theta_R}
    e^{m_t}
    \label{eq:taylor}
\end{equation}
%
where $\bar{r}_B$ is the steady-state nominal interest rate, $\pi^*$ is the
inflation target, $\bar{Y}$ is steady-state output, $\theta_R \in [0,1)$
governs interest rate smoothing, $\theta_\pi > 1$ satisfies the Taylor
principle, and $\theta_Y \geq 0$ governs the response to output fluctuations.

\textbf{Technology shock:}
\begin{equation}
    \log Z_t = \rho_Z \log Z_{t-1} + \eta_t^Z,
    \qquad \eta_t^Z \sim \mathcal{N}(0,\,\sigma_Z^2)
    \label{eq:tfp}
\end{equation}

\textbf{Monetary policy shock:}
\begin{equation}
    m_t = \rho_m\,m_{t-1} + \eta_t^m,
    \qquad \eta_t^m \sim \mathcal{N}(0,\,\sigma_m^2)
    \label{eq:mp_shock}
\end{equation}

% ---------------------------------------------------------------
\section{Model Summary}

Table~\ref{tab:variables} lists all endogenous variables and the equations that
determine them. The baseline model has \textbf{16 endogenous variables} and
\textbf{16 equations}, plus two exogenous shock processes.

\begin{table}[h!]
\centering
\caption{Endogenous variables and determining equations}
\label{tab:variables}
\begin{tabular}{lll}
\toprule
\textbf{Variable} & \textbf{Description} & \textbf{Equation} \\
\midrule
$C_t$               & Consumption                       & CIA (\ref{eq:CIA}) \\
$L_t$               & Hours worked                      & Labor supply (\ref{eq:labor_supply}) \\
$K_{t+1}$           & Capital stock                     & Capital LOM (\ref{eq:capital}) \\
$X_t$               & Investment                        & Investment FOC (\ref{eq:investment_foc}) \\
$M_t$               & Money holdings                    & Money market (\ref{eq:money_market}) \\
$\tilde{\mu}_t$     & CIA multiplier                    & Multiplier split (\ref{eq:mult_split}) \\
$\tilde{\lambda}_t$ & Budget multiplier                 & Multiplier split (\ref{eq:mult_split}) \\
$\tilde{\Theta}_t$  & Shadow price of capital (Tobin's $Q$) & Capital Euler (\ref{eq:capital_euler}) \\
$w_t$               & Real wage                         & Labor demand (\ref{eq:labor_demand}) \\
$r_t$               & Real rental rate of capital       & Capital demand (\ref{eq:capital_demand}) \\
$mc_t$              & Real marginal cost                & Factor demands (\ref{eq:labor_demand})--(\ref{eq:capital_demand}) \\
$Y_t$               & Output                            & Production (\ref{eq:production}) \\
$\pi_t$             & Net inflation rate                & Phillips curve (\ref{eq:phillips}) \\
$r_{B,t}$           & Net nominal interest rate         & Taylor rule (\ref{eq:taylor}) \\
$\tau_t$            & Lump-sum tax rebate               & Fiscal rule (\ref{eq:fiscal}) \\
$P_t$               & Price level                       & Inflation identity (\ref{eq:inflation}) \\
\midrule
$\beta_{t,t+1}$     & Stochastic discount factor        & SDF (\ref{eq:sdf}) \\
\midrule
$Z_t$               & Technology (exogenous)            & TFP process (\ref{eq:tfp}) \\
$m_t$               & Monetary policy shock (exogenous) & MP shock (\ref{eq:mp_shock}) \\
\bottomrule
\end{tabular}
\end{table}

% ---------------------------------------------------------------
\section{Topic Extensions: What Changes}

Table~\ref{tab:extensions} summarises what each topic extension modifies
relative to the baseline, what new variables and equations are introduced, and
which baseline equations are replaced or augmented. Students are responsible for
re-deriving the affected first-order conditions and updating the model summary
table in the appendix of their report accordingly.

\begin{table}[h!]
\centering
\caption{Model extensions by topic}
\label{tab:extensions}
\begin{tabular}{p{1.1cm} p{3.0cm} p{3.8cm} p{2.8cm} p{3.2cm}}
\toprule
\textbf{Topic} & \textbf{Extension} & \textbf{What changes} &
\textbf{New variable(s)} & \textbf{New / modified equation(s)} \\
\midrule
1
  & Variable capital utilization with adjustment costs
  & Production function; capital LOM (\ref{eq:capital}); capital Euler
    (\ref{eq:capital_euler}); investment FOC (\ref{eq:investment_foc})
  & $u_t$; $u_{t-1}$ (predetermined)
  & Utilization FOC; modified capital Euler; modified
    capital LOM \\[0.5em]
2
  & Consumption habit formation (scaled)
  & Utility function; money Euler (\ref{eq:money_euler}); labor supply
    (\ref{eq:labor_supply})
  & $C_{t-1}$ (predetermined)
  & Modified money Euler; modified labor supply \\[0.5em]
3
  & Labor adjustment costs
  & Production technology (\ref{eq:production}); labor demand
    (\ref{eq:labor_demand})
  & $L_{t-1}$ (predetermined)
  & Dynamic labor demand equation \\[0.5em]
4
  & Durable consumption goods
  & Utility function; budget constraint (\ref{eq:budget}); resource
    constraint (\ref{eq:resource})
  & $D_t$ (predetermined); $X_{D,t}$
  & Durable LOM; FOC for durable expenditure \\[0.5em]
5
  & Price indexation
  & Production technology (\ref{eq:production}); Phillips curve
    (\ref{eq:phillips})
  & $\pi_{t-1}$ (predetermined)
  & Modified Phillips curve \\
\bottomrule
\end{tabular}
\end{table}

\end{document}
